%%
%% citemsp-doc.tex — Documentation for the citemsp package
%%
%% Copyright (c) 2025--2026 Apostolos Tsampodimos
%% LPPL 1.3c — see LICENSE for details
%%
\documentclass[11pt, a4paper]{article}

\usepackage{lmodern}
\usepackage{microtype}
\usepackage[margin=2.5cm]{geometry}
\usepackage[dvipsnames, table]{xcolor}
\definecolor{linkblue}{rgb}{0.0, 0.2, 0.6}
\definecolor{linkred}{rgb}{0.7, 0.13, 0.13}
\usepackage{amssymb}
\usepackage[style=numeric, backend=biber]{biblatex}
\usepackage{citemsp}
\usepackage[colorlinks=true, citecolor=linkblue, linkcolor=linkred, urlcolor=linkblue]{hyperref}
\usepackage{tcolorbox}
\usepackage{authblk}
\usepackage{graphicx}
\usepackage{array}
\usepackage{booktabs}
\usepackage{orcidlink}
\usepackage{fontawesome5}
\usepackage{enumitem}

% ── Helpers ───────────────────────────────────────────────────────────────────
\newcommand{\githublink}[1]{\href{https://github.com/ApostolosTsabodimos/citemsp-Latex-package}{\faGithub}}
\newcommand{\pkg}[1]{\texttt{#1}}
\newcommand{\cmd}[1]{\texttt{\textbackslash #1}}

% ── Example bibliography ─────────────────────────────────────────────────────
\begin{filecontents}[force]{refs-citemsp.bib}
@article{upper1974,
  author  = {Upper, Dennis},
  title   = {The Unsuccessful Self-Treatment of a Case of ``Writer's Block''},
  journal = {Journal of Applied Behavior Analysis},
  volume  = {7},
  number  = {3},
  pages   = {497},
  year    = {1974},
  doi     = {10.1901/jaba.1974.7-497a}
}
@article{pennycook2015,
  author  = {Pennycook, Gordon and Cheyne, James Allan and Barr, Nathaniel
             and Koehler, Derek J. and Fugelsang, Jonathan A.},
  title   = {On the Reception and Detection of Pseudo-Profound Bullshit},
  journal = {Judgment and Decision Making},
  volume  = {10},
  number  = {6},
  pages   = {549--563},
  year    = {2015},
  doi     = {10.1017/S1930297500006999}
}
@article{smith2003,
  author  = {Smith, Gordon C. S. and Pell, Jill P.},
  title   = {Parachute Use to Prevent Death and Major Trauma Related to
             Gravitational Challenge: Systematic Review of Randomised
             Controlled Trials},
  journal = {BMJ},
  volume  = {327},
  pages   = {1459--1461},
  year    = {2003},
  doi     = {10.1136/bmj.327.7429.1459}
}
@article{meyerrochow2003,
  author  = {Meyer-Rochow, Victor Benno and Gal, J\'ozsef},
  title   = {Pressures Produced When Penguins Pooh --- Calculations on
             Avian Defaecation},
  journal = {Polar Biology},
  volume  = {27},
  pages   = {56--58},
  year    = {2003},
  doi     = {10.1007/s00300-003-0563-3}
}
@article{wang2020,
  author  = {Wang, Liming and Sofer, Zden\v{e}k and Pumera, Martin},
  title   = {Will Any Crap We Put into Graphene Increase Its
             Electrocatalytic Effect?},
  journal = {ACS Nano},
  volume  = {14},
  number  = {1},
  pages   = {21--25},
  year    = {2020},
  doi     = {10.1021/acsnano.9b00184}
}
@article{cox1979,
  author  = {Cox, David A. and Zucker, Steven},
  title   = {Intersection Numbers of Sections of Elliptic Surfaces},
  journal = {Inventiones Mathematicae},
  volume  = {53},
  pages   = {1--44},
  year    = {1979},
  doi     = {10.1007/BF01403189}
}
@article{zongker2006,
  author  = {Zongker, Doug},
  title   = {Chicken Chicken Chicken: Chicken Chicken},
  journal = {Annals of Improbable Research},
  volume  = {12},
  number  = {5},
  pages   = {16--21},
  year    = {2006},
  url     = {https://isotropic.org/papers/chicken.pdf}
}
@article{ghirlanda2002,
  author  = {Ghirlanda, Stefano and Jansson, Liselotte and Enquist, Magnus},
  title   = {Chickens Prefer Beautiful Humans},
  journal = {Human Nature},
  volume  = {13},
  number  = {3},
  pages   = {383--389},
  year    = {2002},
  doi     = {10.1007/s12110-002-1021-6}
}
@article{matthews1995,
  author  = {Matthews, Robert A. J.},
  title   = {Tumbling Toast, {M}urphy's Law and the Fundamental Constants},
  journal = {European Journal of Physics},
  volume  = {16},
  number  = {4},
  pages   = {172--176},
  year    = {1995},
  doi     = {10.1088/0143-0807/16/4/005}
}
@article{camposarceiz2009,
  author  = {Campos-Arceiz, Ahimsa},
  title   = {Shit Happens (to be Useful)! {U}se of Elephant Dung as
             Habitat by Amphibians},
  journal = {Biotropica},
  volume  = {41},
  number  = {4},
  pages   = {406--407},
  year    = {2009},
  doi     = {10.1111/j.1744-7429.2009.00525.x}
}
@article{yap2022,
  author  = {Yap, Te Faye and Liu, Zhen and Rajappan, Anoop and
             Shimokusu, Trevor J. and Preston, Daniel J.},
  title   = {Necrobotics: Biotic Materials as Ready-to-Use Actuators},
  journal = {Advanced Science},
  volume  = {9},
  pages   = {2201174},
  year    = {2022},
  doi     = {10.1002/advs.202201174}
}
@article{lazebnik2002,
  author  = {Lazebnik, Yuri},
  title   = {Can a Biologist Fix a Radio? --- Or, What {I} Learned While
             Studying Apoptosis},
  journal = {Cancer Cell},
  volume  = {2},
  number  = {3},
  pages   = {179--182},
  year    = {2002},
  doi     = {10.1016/S1535-6108(02)00133-2}
}
@article{mayer2012,
  author  = {Mayer, Hans C. and Krechetnikov, Rouslan},
  title   = {Walking with Coffee: Why Does It Spill?},
  journal = {Physical Review E},
  volume  = {85},
  pages   = {046117},
  year    = {2012},
  doi     = {10.1103/PhysRevE.85.046117}
}
\end{filecontents}
\addbibresource{refs-citemsp.bib}

% ── Title ─────────────────────────────────────────────────────────────────────
\title{\textbf{\pkg{citemsp} Package}\\[4pt]
  \large Per-Key Citation Locators\\[4pt]
  \normalsize v2.3 --- \today}

\author[1]{Apostolos Tsampodimos\,\,\githublink{}\,\,\orcidlink{0009-0001-0230-5647}}
\author[2]{Nathaniel Sherrill}
\author[2]{Agnese Mariotti}

\affil[1]{Center for Physical Sciences and Technology\\
Saul\.etekio~3, 10257 Vilnius, Lithuania\\
\texttt{apostolos.tsampodimos@ftmc.lt}}
\affil[2]{Institute for Theoretical Physics, Leibniz Universit\"at Hannover\\
Appelstra\ss e~2, 30167 Hannover, Germany}

\date{}

\begin{document}

\maketitle

\begin{abstract}
\noindent
The \pkg{citemsp} package provides a single user-facing command, \cmd{citemsp}, that attaches section~(\S), paragraph~(\P), and prefix-based locators directly to numeric citation labels. It works on top of either \pkg{biblatex} or \pkg{natbib}, auto-detecting which backend is loaded. An extensible prefix registry provides, currently, ten built-in types and users can register custom prefixes with \cmd{citemspprefix}. The package parses a slash-delimited syntax, rendering each locator as a superscript/subscript pair on the corresponding citation number, producing output such as~\citemsp{smith2003/3.2/5, wang2020/p520, meyerrochow2003/1/e4.3}. Plain entries (no slashes) are batched and passed to the backend's native \cmd{cite}, so range compression (\pkg{numeric-comp}, \texttt{sort\&compress}) works naturally.
\end{abstract}

%\tableofcontents

\newpage

% ══════════════════════════════════════════════════════════════════════════════
\section{Motivation}
% ══════════════════════════════════════════════════════════════════════════════

The \pkg{biblatex} package offers a locator interface through its postnote
field.  A user can write \verb|\autocite[§3.2]{smith2003}| to append a
free-form locator after the citation number.  For multi-key citations,
\verb|\autocites| accepts per-key pre- and post-notes:
%
\begin{tcolorbox}[colback=gray!5, colframe=gray!40, boxrule=0.5pt, arc=2pt,
  top=6pt, bottom=6pt, left=10pt, right=10pt]
\begin{verbatim}
\autocites[§3.2]{smith2003}[¶2]{meyerrochow2003}
\end{verbatim}
\end{tcolorbox}
%
\noindent This approach has two limitations when writing a scientific paper:

\begin{enumerate}[leftmargin=2em]
  \item \textbf{Verbose syntax.}  Citing three sources with locators
    requires nested optional arguments, which is cumbersome in source files.
  \item \textbf{No typed locators.}  There is no compact notation for
    referencing specific equations, figures, definitions, or theorems within
    cited works.
\end{enumerate}

\noindent The \pkg{citemsp} package renders locators as a
superscript/subscript pair attached to the citation number, supports ten
built-in prefix types via an extensible registry, and accepts all keys with
their locators in a single comma-separated argument.  It supports both
\pkg{biblatex} and \pkg{natbib} backends.


% ══════════════════════════════════════════════════════════════════════════════
\section{Installation}
% ══════════════════════════════════════════════════════════════════════════════

\subsection{Per-project}

Copy \texttt{citemsp.sty} into the same directory as your main
\texttt{.tex} file.

\subsection{System-wide (local texmf)}

\begin{tcolorbox}[colback=gray!5, colframe=gray!40, boxrule=0.5pt, arc=2pt,
  top=6pt, bottom=6pt, left=10pt, right=10pt]
\begin{verbatim}
mkdir -p ~/texmf/tex/latex/citemsp
cp citemsp.sty ~/texmf/tex/latex/citemsp/
texhash ~/texmf
\end{verbatim}
\end{tcolorbox}


% ══════════════════════════════════════════════════════════════════════════════
\section{Usage}
% ══════════════════════════════════════════════════════════════════════════════

\subsection{Loading with biblatex}

The package must be loaded \emph{after} \pkg{biblatex}:
%
\begin{tcolorbox}[colback=gray!5, colframe=gray!40, boxrule=0.5pt, arc=2pt,
  top=6pt, bottom=6pt, left=10pt, right=10pt]
\begin{verbatim}
\usepackage[style=numeric, backend=biber]{biblatex}
\addbibresource{refs.bib}
\usepackage{citemsp}
\end{verbatim}
\end{tcolorbox}

\subsection{Loading with natbib}

\begin{tcolorbox}[colback=gray!5, colframe=gray!40, boxrule=0.5pt, arc=2pt,
  top=6pt, bottom=6pt, left=10pt, right=10pt]
\begin{verbatim}
\usepackage[numbers]{natbib}   % omit if your class loads it
\usepackage{citemsp}
\end{verbatim}
\end{tcolorbox}
%
\noindent Classes that auto-load \pkg{natbib} (e.g., \pkg{revtex4-2},
\pkg{aastex631}) require no explicit \verb|\usepackage{natbib}|.
Loading \pkg{citemsp} before either backend produces a package error.

\subsection{Syntax}

The \cmd{citemsp} command accepts a comma-separated list of entries.
Each entry has the form \texttt{key}, \texttt{key/loc1}, or
\texttt{key/loc1/loc2}.  By default, the first slot renders as
\S\ (section) and the second as \P\ (paragraph).  A single-letter prefix
changes the locator type:

\begin{tcolorbox}[colback=gray!5, colframe=gray!40, boxrule=0.5pt, arc=2pt,
  top=6pt, bottom=6pt, left=10pt, right=10pt]
\begin{tabular*}{\dimexpr\linewidth-0pt}{@{}l@{\extracolsep{\fill}}c@{\quad}l@{}}
\texttt{\textbackslash citemsp\{key\}}
  & $\longrightarrow$ & \citemsp{smith2003}\\[3pt]
\texttt{\textbackslash citemsp\{key/section\}}
  & $\longrightarrow$ & \citemsp{smith2003/6.2}\\[3pt]
\texttt{\textbackslash citemsp\{key/section/paragraph\}}
  & $\longrightarrow$ & \citemsp{smith2003/6.2/3}\\[3pt]
\texttt{\textbackslash citemsp\{key//paragraph\}}
  & $\longrightarrow$ & \citemsp{smith2003//3}\\[3pt]
\texttt{\textbackslash citemsp\{key/section/e4.3\}}
  & $\longrightarrow$ & \citemsp{smith2003/6.2/e4.3}
\end{tabular*}
\end{tcolorbox}

\noindent Note the subscript-only form \texttt{key//par}: an empty first
slot with a double slash gives a paragraph (or prefixed) subscript without
a section superscript.


\subsection{Prefix registry}

All ten built-in prefixes are listed below. Note that no prefix defaults to \S\ (section) in the superscript slot and \P\ (paragraph) in the subscript slot. A prefix is a single letter at the start of a locator slot that selects a symbol:

\begin{table}[h]
\centering
\rowcolors{2}{blue!3}{white}
\begin{tabular}{c l l c c c}
\toprule
\rowcolor{blue!15}
\textbf{Prefix} & \textbf{Locator} & \textbf{Render} & \textbf{In super slot} & \textbf{In sub slot} & \textbf{Both slots} \\
\midrule
\texttt{A} & Appendix    & App.          & \citemsp{smith2003/A2}        & \citemsp{smith2003//A2}         & \citemsp{smith2003/3.2/A2} \\
\texttt{C} & Corollary   & Cor.          & \citemsp{smith2003/C3}        & \citemsp{smith2003//C3}      & \citemsp{smith2003/4.1/C3} \\
\texttt{L} & Lemma       & Lem.          & \citemsp{smith2003/L1}        & \citemsp{smith2003//L1}      & \citemsp{smith2003/3.2/L1} \\
\texttt{T} & Theorem     & Th.           & \citemsp{smith2003/T4.2}      & \citemsp{smith2003//T4.2}    & \citemsp{smith2003/6.1/T4.2} \\
\texttt{d} & Definition  & $\triangleq$  & \citemsp{smith2003/d5}        & \citemsp{smith2003//d5}      & \citemsp{smith2003/3.2/d5} \\
\texttt{e} & Equation    & eq.           & \citemsp{smith2003/e4.3}      & \citemsp{smith2003//e4.3}    & \citemsp{smith2003/6.2/e4.3} \\
\texttt{f} & Figure      & fig.          & \citemsp{smith2003/f3}        & \citemsp{smith2003//f3}      & \citemsp{smith2003/6.1/f3} \\
\texttt{n} & Footnote    & $\dagger$     & \citemsp{meyerrochow2003/n7}     & \citemsp{meyerrochow2003//n7}     & \citemsp{meyerrochow2003/1/n7} \\
\texttt{p} & Page        & pp.           & \citemsp{smith2003/p42}       & \citemsp{smith2003//p42}     & \citemsp{smith2003/3.2/p42} \\
\texttt{t} & Table       & $\boxplus$    & \citemsp{wang2020/t2}         & \citemsp{wang2020//t2}       & \citemsp{wang2020/8.4/t2} \\
\bottomrule
\end{tabular}
\caption{Built-in prefix registry.}
\label{tab:prefixes}
\end{table}


\noindent Users can register custom prefixes with \cmd{citemspprefix}.
Place the command in the preamble (after \verb|\usepackage{citemsp}|) so
it is available throughout the document:

\begin{tcolorbox}[colback=gray!5, colframe=gray!40, boxrule=0.5pt, arc=2pt,
  top=6pt, bottom=6pt, left=10pt, right=10pt]
\begin{verbatim}
\usepackage{citemsp}
\citemspprefix{R}{Rem.}         "R" renders as "Rem."
\citemspprefix{w}{w.}           "w" renders as "w."
\end{verbatim}
\end{tcolorbox}

\citemspprefix{R}{Rem.}%
\noindent For example, after registering \texttt{R} for ``Remark'':
\verb|\citemsp{cox1979/12.4/R3}| $\longrightarrow$ \citemsp{cox1979/12.4/R3},\enspace
and after registering \texttt{w} for ``widget'':
\verb|\citemsp{zongker2006/3/w5}| $\longrightarrow$ \citemsp{zongker2006/3/w5}.\enspace
The prefix letter must not collide with an existing built-in prefix.

\subsection{Configuration}

Two user-configurable parameters control the locator appearance:
\begin{itemize}[leftmargin=2em]
  \item \cmd{citemspscale} (default~0.35) --- controls the locator glyph
    size via \cmd{scalebox} from \pkg{graphicx}.
  \item \cmd{citemspraiseoffset} (default~1.5pt) --- vertical offset
    applied via \cmd{raisebox} to superscript locators.
\end{itemize}

\noindent Redefine them after loading:
\begin{tcolorbox}[colback=gray!5, colframe=gray!40, boxrule=0.5pt, arc=2pt,
  top=6pt, bottom=6pt, left=10pt, right=10pt]
\begin{verbatim}
\renewcommand{\citemspscale}{0.4}
\renewcommand{\citemspraiseoffset}{2pt}
\end{verbatim}
\end{tcolorbox}


% ══════════════════════════════════════════════════════════════════════════════
\section{Comparison with Standard Biblatex}
% ══════════════════════════════════════════════════════════════════════════════

\begin{table}[ht]
\centering
\rowcolors{2}{blue!3}{white}
\begin{tabular}{>{\raggedright\arraybackslash}p{0.46\textwidth}
                >{\raggedright\arraybackslash}p{0.46\textwidth}}
\toprule
\rowcolor{blue!15}
\textbf{Standard biblatex} & \textbf{citemsp} \\
\midrule
\verb|\autocite[§3.2]{smith2003}| \newline
$\rightarrow$ [1, \S3.2]
&
\verb|\citemsp{smith2003/3.2}| \newline
$\rightarrow$ \citemsp{smith2003/3.2}
\\[6pt]
\verb|\autocites[§3.2]{smith2003}| \verb|[eq.~4.3]{meyerrochow2003}| \newline
$\rightarrow$ [1, \S3.2][2, eq.~4.3]
&
\verb|\citemsp{smith2003/3.2,| \verb|meyerrochow2003/e4.3}| \newline
$\rightarrow$ \citemsp{smith2003/3.2, meyerrochow2003/e4.3}
\\[6pt]
\verb|\autocites[§3.2, ¶5]{smith2003}| \newline
$\rightarrow$ [1, \S3.2, \P5]
&
\verb|\citemsp{smith2003/3.2/5}| \newline
$\rightarrow$ \citemsp{smith2003/3.2/5}
\\[6pt]
No built-in typed locator syntax; user must type symbols manually
&
Ten built-in prefixes; extensible via \cmd{citemspprefix}
\\[6pt]
Range compression via \texttt{numeric-comp}
&
Plain entries (no~\texttt{/}) compressed via the backend's native
\cmd{cite}, so compression works identically
\\
\bottomrule
\end{tabular}
\caption{Feature comparison between standard \pkg{biblatex} locators and
  \pkg{citemsp}.}
\label{tab:comparison}
\end{table}


% ══════════════════════════════════════════════════════════════════════════════
\section{Examples}
% ══════════════════════════════════════════════════════════════════════════════

\subsection{Single source with section and paragraph}

\begin{tcolorbox}[colback=blue!3, colframe=blue!30, boxrule=0.5pt, arc=2pt,
  top=6pt, bottom=6pt, left=10pt, right=10pt]
\begin{verbatim}
The lack of randomised controlled trials on
parachutes~\citemsp{smith2003/6.1/2} remains a gap in
evidence-based medicine.
\end{verbatim}
\end{tcolorbox}

\noindent\textbf{Output:} The lack of randomised controlled trials on
parachutes~\citemsp{smith2003/6.1/2} remains a gap in evidence-based
medicine.


\subsection{Section-only locator}

\begin{tcolorbox}[colback=blue!3, colframe=blue!30, boxrule=0.5pt, arc=2pt,
  top=6pt, bottom=6pt, left=10pt, right=10pt]
\begin{verbatim}
Chickens were found to prefer beautiful
humans~\citemsp{ghirlanda2002/2.3}, raising questions about
peer review in poultry science.
\end{verbatim}
\end{tcolorbox}

\noindent\textbf{Output:} Chickens were found to prefer beautiful
humans~\citemsp{ghirlanda2002/2.3}, raising questions about peer review
in poultry science.


\subsection{Subscript-only locator}

An empty first slot (double slash) gives a subscript without a superscript:

\begin{tcolorbox}[colback=blue!3, colframe=blue!30, boxrule=0.5pt, arc=2pt,
  top=6pt, bottom=6pt, left=10pt, right=10pt]
\begin{verbatim}
\citemsp{smith2003//5}              % paragraph 5, no section
\citemsp{meyerrochow2003//e2.1}        % equation 2.1, no section
\end{verbatim}
\end{tcolorbox}

\noindent\textbf{Output:}
\citemsp{smith2003//5},\enspace
\citemsp{meyerrochow2003//e2.1}.


\subsection{Multiple sources with mixed locators}

\begin{tcolorbox}[colback=blue!3, colframe=blue!30, boxrule=0.5pt, arc=2pt,
  top=6pt, bottom=6pt, left=10pt, right=10pt]
\begin{verbatim}
Writer's block~\citemsp{upper1974, smith2003/4,
  meyerrochow2003/1/2} has been studied from behavioural,
medical, and ornithological perspectives.
\end{verbatim}
\end{tcolorbox}

\noindent\textbf{Output:} Writer's
block~\citemsp{upper1974, smith2003/4, meyerrochow2003/1/2} has been
studied from behavioural, medical, and ornithological perspectives.


\subsection{Range compression}

Plain entries (no slashes) are batched and passed to the backend's native
\cmd{cite}, so range compression works automatically when using
\texttt{numeric-comp} (biblatex) or \texttt{sort\&compress} (natbib).
Entries \emph{with} locators are rendered individually and interleaved
with the compressed group.  The examples below illustrate different
combinations.

\medskip
\noindent\textbf{1.\enspace All plain (fully compressed):}

\begin{tcolorbox}[colback=blue!3, colframe=blue!30, boxrule=0.5pt, arc=2pt,
  top=6pt, bottom=6pt, left=10pt, right=10pt]
\begin{verbatim}
\citemsp{upper1974, pennycook2015, smith2003, meyerrochow2003,
  wang2020}
\end{verbatim}
\end{tcolorbox}

\noindent\textbf{Output:}
\citemsp{upper1974, pennycook2015, smith2003, meyerrochow2003, wang2020}

\medskip
\noindent\textbf{2.\enspace Plain entries mixed with locator entries:}

\begin{tcolorbox}[colback=blue!3, colframe=blue!30, boxrule=0.5pt, arc=2pt,
  top=6pt, bottom=6pt, left=10pt, right=10pt]
\begin{verbatim}
\citemsp{upper1974, pennycook2015, smith2003/4.1,
  meyerrochow2003, wang2020}
\end{verbatim}
\end{tcolorbox}

\noindent\textbf{Output:}
\citemsp{upper1974, pennycook2015, smith2003/4.1, meyerrochow2003, wang2020}

\smallskip\noindent
The plain entries compress into a range while the locator entry renders
separately.

\medskip
\noindent\textbf{3.\enspace Multiple locator entries among plain ones:}

\begin{tcolorbox}[colback=blue!3, colframe=blue!30, boxrule=0.5pt, arc=2pt,
  top=6pt, bottom=6pt, left=10pt, right=10pt]
\begin{verbatim}
\citemsp{upper1974, pennycook2015/2.3, smith2003,
  meyerrochow2003/e14, wang2020}
\end{verbatim}
\end{tcolorbox}

\noindent\textbf{Output:}
\citemsp{upper1974, pennycook2015/2.3, smith2003, meyerrochow2003/e14, wang2020}

\medskip
\noindent\textbf{4.\enspace Every entry has a locator (no compression):}

\begin{tcolorbox}[colback=blue!3, colframe=blue!30, boxrule=0.5pt, arc=2pt,
  top=6pt, bottom=6pt, left=10pt, right=10pt]
\begin{verbatim}
\citemsp{upper1974/1, smith2003/4.1/2,
  meyerrochow2003/e14, wang2020/p21}
\end{verbatim}
\end{tcolorbox}

\noindent\textbf{Output:}
\citemsp{upper1974/1, smith2003/4.1/2, meyerrochow2003/e14, wang2020/p21}

\smallskip\noindent
When every entry carries a locator, each is rendered individually and no
range compression occurs.

\medskip
\noindent\textbf{5.\enspace Locator entries with super and subscript among plain ones:}

\begin{tcolorbox}[colback=blue!3, colframe=blue!30, boxrule=0.5pt, arc=2pt,
  top=6pt, bottom=6pt, left=10pt, right=10pt]
\begin{verbatim}
\citemsp{upper1974, pennycook2015, smith2003/6.1/2,
  meyerrochow2003, wang2020, cox1979/12.4/T5}
\end{verbatim}
\end{tcolorbox}

\noindent\textbf{Output:}
\citemsp{upper1974, pennycook2015, smith2003/6.1/2, meyerrochow2003, wang2020, cox1979/12.4/T5}


\subsection{Equation and figure locators}

\begin{tcolorbox}[colback=blue!3, colframe=blue!30, boxrule=0.5pt, arc=2pt,
  top=6pt, bottom=6pt, left=10pt, right=10pt]
\begin{verbatim}
The sloshing dynamics of coffee~\citemsp{mayer2012/3.2/e3.31}
and the critical spill diagram~\citemsp{mayer2012/4.6/f4.1}
explain why your laptop is always sticky.
\end{verbatim}
\end{tcolorbox}

\noindent\textbf{Output:} The sloshing dynamics of
coffee~\citemsp{mayer2012/3.2/e3.31} and the critical spill
diagram~\citemsp{mayer2012/4.6/f4.1} explain why your laptop is
always sticky.


\subsection{Theorem, lemma, and corollary locators}

\begin{tcolorbox}[colback=blue!3, colframe=blue!30, boxrule=0.5pt, arc=2pt,
  top=6pt, bottom=6pt, left=10pt, right=10pt]
\begin{verbatim}
The Cox--Zucker machine~\citemsp{cox1979/12.4/T5} computes
intersection numbers; see Lemma~\citemsp{cox1979/12.3/L2}
and Corollary~\citemsp{cox1979/12.4/C3}.
\end{verbatim}
\end{tcolorbox}

\noindent\textbf{Output:} The Cox--Zucker
machine~\citemsp{cox1979/12.4/T5} computes intersection numbers; see
Lemma~\citemsp{cox1979/12.3/L2} and
Corollary~\citemsp{cox1979/12.4/C3}.


\subsection{Definition and page locators}

\begin{tcolorbox}[colback=blue!3, colframe=blue!30, boxrule=0.5pt, arc=2pt,
  top=6pt, bottom=6pt, left=10pt, right=10pt]
\begin{verbatim}
The definition of a necrobotic
actuator~\citemsp{yap2022/2.3/d2.1} involves reanimating
dead spiders; see also~\citemsp{yap2022/p37}.
\end{verbatim}
\end{tcolorbox}

\noindent\textbf{Output:} The definition of a necrobotic
actuator~\citemsp{yap2022/2.3/d2.1} involves reanimating dead
spiders; see also~\citemsp{yap2022/p37}.


\subsection{Appendix and footnote locators}

\begin{tcolorbox}[colback=blue!3, colframe=blue!30, boxrule=0.5pt, arc=2pt,
  top=6pt, bottom=6pt, left=10pt, right=10pt]
\begin{verbatim}
The theorem that biologists cannot fix a
radio~\citemsp{lazebnik2002/T11} is proven in
Appendix~\citemsp{lazebnik2002/A2}; the apology to
electrical engineers is in
footnote~\citemsp{lazebnik2002/3/n4}.
\end{verbatim}
\end{tcolorbox}

\noindent\textbf{Output:} The theorem that biologists cannot fix a
radio~\citemsp{lazebnik2002/T11} is proven in
Appendix~\citemsp{lazebnik2002/A2}; the apology to electrical engineers
is in footnote~\citemsp{lazebnik2002/3/n4}.


\subsection{Table locator and mixed types in one call}

\begin{tcolorbox}[colback=blue!3, colframe=blue!30, boxrule=0.5pt, arc=2pt,
  top=6pt, bottom=6pt, left=10pt, right=10pt]
\begin{verbatim}
See~\citemsp{smith2003/6.1/2, matthews1995/e6.55,
  wang2020/8.4/t2, yap2022/p170} for parachutes, toast,
graphene, and dead spiders respectively.
\end{verbatim}
\end{tcolorbox}

\noindent\textbf{Output:}
See~\citemsp{smith2003/6.1/2, matthews1995/e6.55, wang2020/8.4/t2,
yap2022/p170} for parachutes, toast, graphene, and dead spiders
respectively.


\subsection{Plain citation as drop-in replacement}

Without slashes, \cmd{citemsp} behaves identically to \cmd{cite}:

\begin{tcolorbox}[colback=blue!3, colframe=blue!30, boxrule=0.5pt, arc=2pt,
  top=6pt, bottom=6pt, left=10pt, right=10pt]
\begin{verbatim}
Elephant dung turns out to be useful
habitat~\citemsp{camposarceiz2009}, which is more than can be
said for most conference proceedings.
\end{verbatim}
\end{tcolorbox}

\noindent\textbf{Output:} Elephant dung turns out to be useful
habitat~\citemsp{camposarceiz2009}, which is more than can be said for most
conference proceedings.


\subsection{Adjacent to inline math}

\begin{tcolorbox}[colback=blue!3, colframe=blue!30, boxrule=0.5pt, arc=2pt,
  top=6pt, bottom=6pt, left=10pt, right=10pt]
\begin{verbatim}
The penguin defaecation pressure $P \approx 60\,
\mathrm{kPa}$~\citemsp{meyerrochow2003/e14, smith2003/4.3}
exceeds most expectations.
\end{verbatim}
\end{tcolorbox}

\noindent\textbf{Output:} The penguin defaecation pressure
$P \approx 60\,\mathrm{kPa}$~\citemsp{meyerrochow2003/e14, smith2003/4.3}
exceeds most expectations.


\subsection{Font contexts}

The package works correctly in all standard \LaTeX\ font contexts.
For italic and slanted shapes, a small kern is automatically inserted to
prevent the locator from colliding with the slanted closing bracket.
Bold, bold-italic, small caps, and upright text work without modification.

\begin{tcolorbox}[colback=blue!3, colframe=blue!30, boxrule=0.5pt, arc=2pt,
  top=6pt, bottom=6pt, left=10pt, right=10pt]
Upright:~\citemsp{smith2003/6.1/2}\\[3pt]
\textit{Italic:~\citemsp{smith2003/6.1/2}}\\[3pt]
\textsl{Slanted:~\citemsp{smith2003/6.1/2}}\\[3pt]
\textbf{Bold:~\citemsp{smith2003/6.1/2}}\\[3pt]
\textbf{\textit{Bold italic:~\citemsp{smith2003/6.1/2}}}\\[3pt]
\textsc{Small caps:~\citemsp{smith2003/6.1/2}}\\[3pt]
\emph{Emph (italic):~\citemsp{smith2003/6.1/2}
  --- \emph{nested emph (upright):~\citemsp{smith2003/6.1/2}}}
\end{tcolorbox}


\subsection{Font sizes}

Locator glyphs scale proportionally at different font sizes:

\begin{tcolorbox}[colback=blue!3, colframe=blue!30, boxrule=0.5pt, arc=2pt,
  top=6pt, bottom=6pt, left=10pt, right=10pt]
{\small Small:~\citemsp{ghirlanda2002/2.3/e2.5}}\\
{\normalsize Normal:~\citemsp{ghirlanda2002/2.3/e2.5}}\\
{\large Large:~\citemsp{ghirlanda2002/2.3/e2.5}}
\end{tcolorbox}


\subsection{Custom prefixes}

Place \cmd{citemspprefix} in the preamble, after loading \pkg{citemsp}:

\citemspprefix{w}{w.}%
\citemspprefix{H}{Hyp.}%
\begin{tcolorbox}[colback=blue!3, colframe=blue!30, boxrule=0.5pt, arc=2pt,
  top=6pt, bottom=6pt, left=10pt, right=10pt]
\begin{verbatim}
% in the preamble:
\citemspprefix{w}{w.}       % "w" renders as "w."
\citemspprefix{H}{Hyp.}     % "H" renders as "Hyp."

% in the document body:
\citemsp{zongker2006/3/w5}
\citemsp{cox1979/12.4/H2}
\end{verbatim}
\end{tcolorbox}

\noindent\textbf{Output:}
\citemsp{zongker2006/3/w5},\enspace
\citemsp{cox1979/12.4/H2}


% ══════════════════════════════════════════════════════════════════════════════
\section{Natbib Backend}
% ══════════════════════════════════════════════════════════════════════════════

Since v2.2, \pkg{citemsp} auto-detects whether \pkg{biblatex} or
\pkg{natbib} is loaded and selects the matching code path.  No package
option is needed.

\subsection{Compatible classes}

Many journal classes auto-load \pkg{natbib}.  With these, simply loading
\pkg{citemsp} after the class is sufficient:

\begin{table}[h]
\centering
\rowcolors{2}{blue!3}{white}
\begin{tabular}{l l l}
\toprule
\rowcolor{blue!15}
\textbf{Class} & \textbf{Backend} & \textbf{Notes} \\
\midrule
\texttt{article}, \texttt{report}, \texttt{book}, KOMA
  & either & user's choice \\
\texttt{revtex4-2} (APS)
  & natbib (auto) & class loads natbib \\
\texttt{aastex631} (AAS)
  & natbib (auto) & class loads natbib \\
\texttt{elsarticle} (Elsevier)
  & natbib (manual) & add \verb|\usepackage[numbers]{natbib}| \\
\texttt{IEEEtran} natbib mode
  & natbib (manual) & add \verb|\usepackage[numbers]{natbib}| \\
\bottomrule
\end{tabular}
\caption{Document class compatibility.}
\label{tab:classes}
\end{table}

\subsection{Minimal natbib example}

\begin{tcolorbox}[colback=gray!5, colframe=gray!40, boxrule=0.5pt, arc=2pt,
  top=6pt, bottom=6pt, left=10pt, right=10pt]
\begin{verbatim}
\documentclass{article}
\usepackage[numbers, sort&compress]{natbib}
\usepackage{citemsp}
\begin{document}
See~\citemsp{key1/3.2/5, key2/e4.1} for details.
\bibliographystyle{unsrt}
\bibliography{refs}
\end{document}
\end{verbatim}
\end{tcolorbox}

\noindent On the natbib path, \pkg{citemsp} uses \cmd{citenum} for the
bare number and wraps it in a \cmd{hyperlink} when \pkg{hyperref} is
loaded.  Range compression via \texttt{sort\&compress} works for plain
entries.


% ══════════════════════════════════════════════════════════════════════════════
\section{Implementation}
% ══════════════════════════════════════════════════════════════════════════════

\begin{enumerate}[leftmargin=2em]
  \item \textbf{Backend detection.}  At load time, the package checks
    whether \pkg{biblatex} or \pkg{natbib} is already loaded.  If neither
    is present, a descriptive error is raised.  If both are loaded
    (itself an error), then \pkg{biblatex} is preferred.

  \item \textbf{Locator scaling.}  \cmd{citemspscale} (default~0.35)
    controls the size of locator glyphs via \cmd{scalebox} from
    \pkg{graphicx}.  \cmd{citemspraiseoffset} (default~1.5pt)
    sets the vertical offset of superscript locators via \cmd{raisebox}.

  \item \textbf{Font shape detection.}  The renderer checks
    \verb|\f@shape| against \texttt{it} (italic) and \texttt{sl}
    (slanted).  In either case, a small math kern is inserted to prevent
    the locator from colliding with the slanted final bracket.

  \item \textbf{Prefix dispatch.}  Each locator value is tested: if its
    first character matches a registered prefix (via
    \verb|\csname @citemsp@pfx@...| lookup), the prefix symbol is
    rendered; otherwise the default \S\ or \P\ is used.

  \item \textbf{Citation rendering.}  On the biblatex path, a custom cite
    command \cmd{citelinksp} (via \verb|\DeclareCiteCommand|) prints the
    hyperlinked number and attaches locators.  On the natbib path,
    \cmd{citenum} provides the bare number, wrapped in
    \cmd{hyperlink} when \pkg{hyperref} is loaded.

  \item \textbf{Entry parsing.}  The user-facing \cmd{citemsp} uses
    \verb|\forcsvlist| (from \pkg{etoolbox}) to iterate over
    comma-separated entries.  Each entry is split at forward slashes
    using a delimited \verb|\def| pattern, extracting the key and up to
    two locator slots.

  \item \textbf{Plain-key batching.}  Entries without slashes are
    accumulated into a buffer and flushed via the backend's native
    \cmd{cite} (with brackets suppressed), so that range compression
    (\texttt{numeric-comp} or \texttt{sort\&compress}) works identically
    to a regular \cmd{cite} call.
\end{enumerate}


% ══════════════════════════════════════════════════════════════════════════════
\section{Limitations}
% ══════════════════════════════════════════════════════════════════════════════

\begin{itemize}[leftmargin=2em]
  \item \textbf{Numeric styles only.}  \pkg{citemsp} attaches locators to numeric citation labels. Other non-numeric styles are currently not supported.
  \item \textbf{Manual paragraph counting.}  Paragraph locators refer to paragraph numbers counted manually from the start of the referenced section.  There is no automatic cross-referencing.
  \item \textbf{\texttt{revtex4-2} with biblatex.}  The combination of \texttt{revtex4-2} and \pkg{biblatex} is incompatible due to conflicts in \verb|\addtocontents| patching between the two packages.  Although not a \pkg{citemsp} issue, users should keep that in mind. Make sure to use the natbib backend with revtex.
\end{itemize}


% ══════════════════════════════════════════════════════════════════════════════
\section{Conventions}
% ══════════════════════════════════════════════════════════════════════════════

\begin{itemize}[leftmargin=2em]
  \item \textbf{Section numbering:} use the source's native format
    (e.g., 3.2.1, IV, A.1).
  \item \textbf{Paragraph counting:} count paragraphs from the start of
    the referenced section.  This is a manual convention; \pkg{citemsp}
    does not verify paragraph numbers against the cited source.
  \item \textbf{Subsection references:} use the full path
    (e.g., \texttt{key/3.2.1}), not a relative number.
  \item \textbf{Plain citations:} \verb|\citemsp{key}| with no slashes
    produces a standard numeric citation, so the command can serve as a
    drop-in replacement for \cmd{cite} if desired.
\end{itemize}


% ══════════════════════════════════════════════════════════════════════════════
\section{Requirements}
% ══════════════════════════════════════════════════════════════════════════════

\begin{itemize}[leftmargin=2em]
  \item \LaTeX2e\ (2020/10/01 or later)
  \item \pkg{biblatex} (any numeric style) \textbf{or} \pkg{natbib}
    (with the \texttt{numbers} option)
  \item \pkg{graphicx}
  \item \pkg{etoolbox}
  \item \pkg{amssymb} --- recommended for best glyph quality
    ($\triangleq$, $\boxplus$).  If not loaded, \pkg{citemsp} provides
    primitive fallback definitions.
\end{itemize}


% ══════════════════════════════════════════════════════════════════════════════
\section{License}
% ══════════════════════════════════════════════════════════════════════════════

This work is released under the
\href{https://www.latex-project.org/lppl/lppl-1-3c/}{LaTeX Project Public
License, version~1.3c} or later.

\newpage

\printbibliography

\end{document}
